\paragraph{Two distinct BE failure diagnoses (post-test).}
In the two deteriorated CF seeds, 105001 and 105008, COMMIT BA on the 3,600
acquired auxiliary queries is 99.510\% and 100.000\%, versus 52.756\%
and 52.094\% on the held-out contexts. Both models visited all auxiliary
queries. Failure to fit acquired labels alone does not explain these cases:
the observed COMMIT gap concerns generalization across contexts, without
identifying its optimization cause or an architectural impossibility.
All 181 final false negatives are at the eligibility boundary
$\kappa=\operatorname{sgn}(a_v-q)=0$ and have no participating neighbor.
They are distinct from the active-neighbor defect in the separate
follow-up below. On eligible silent queries CF makes 81 false acceptances
versus 42 for native; these bits alone are masked by participation.

Seed 105010 has a different defect: both arms complete only 204/1,000 runs,
despite final COMMIT BA 99.451\% native and 99.949\% CF. Their shared
first stage omits participation on 13/118 eligible test roots.
A replay of the original 32-case cohort localizes the first invariant
violation in one selected run: at KW step 7, vertex 31 must
recolor as the bound drops from 21 to 17, but remains at 17
because participation is false. Its proposed value 12 is free;
changing only root participation makes COMMIT accept it. Thus this state
exposes an upstream omission, rather than an incorrect veto of the actual
silent message. These are selected post-test diagnostics, without retuning,
a continued rescue rollout or an estimate of causes across all failed runs
(Appendix~\ref{app:be-failure-diagnostics}).
